\section{Summary of Findings}

\subsection{Summary of Contributions}

\paragraph{1. Integration of StyleFeatureEditor (SFE) Encoder} 
FeatureSAGE integrates StyleFeatureEditor Encoder on the original SAGE framework, enabling improved latent representations that better 
capture distributional characteristics of the data across resolutions.

\paragraph{2. Integration of Feature Editor} 
The Feature Editor module allows fine-grained and semantically consistent modifications by operating directly on intermediate feature representations, enhancing the 
expressiveness of attribute edits.

\paragraph{3. Resolution-Dependent and Dataset-Adaptive Editing} 
FeatureSAGE adapts its editing strategies according to both image resolution and dataset characteristics, ensuring appropriate trade-offs between edit diversity, 
stability, and realism.

\paragraph{4. Comprehensive Benchmarking Against Baselines} 
The framework demonstrates robust performance across structurally diverse datasets (Flowers and Animal Faces) and multiple resolutions, providing a generalizable 
approach to few-shot image generation while balancing perceptual expressiveness and identity preservation.

\subsection{Empirical Outcomes}

\paragraph{Quantitative Performance} 
FeatureSAGE consistently improves distributional realism across datasets and resolutions. Key metrics include:

\begin{itemize}
    \item \textbf{FID improvements:} At $1024 \times 1024$, FeatureSAGE reduces FID by 25.9\% on Flowers (45.60 vs. 61.57 for AGE) and 15.8\% (vs. 54.17 for SAGE), and by 20.2\% on Animal Faces (49.60 vs. 62.13 for AGE) and 18.6\% (vs. 60.92 for SAGE). At $256 \times 256$ on Flowers, reductions are 15.0\% (vs. AGE) and 9.1\% (vs. SAGE).
    \item \textbf{LPIPS performance:} On Flowers at $1024 \times 1024$, FeatureSAGE achieves an LPIPS of 0.4209, which is 7.0\% lower than SAGE (0.4528) and 2.5\% higher than AGE (0.4108), maintaining competitive perceptual diversity. At $256 \times 256$, LPIPS decreases to 0.1309, representing a 37.2\% reduction compared to SAGE (0.2085) and 37.9\% compared to AGE (0.2108). On Animal Faces at $1024 \times 1024$, LPIPS is 0.4811, within 1.2\% of SAGE (0.4867) and 4.4\% of AGE (0.5033).
\end{itemize}

\paragraph{Dataset and Resolution-Adaptive Trends} 
The framework demonstrates resolution-dependent behavior: 
\begin{itemize}
    \item At $1024 \times 1024$ resolution, FeatureSAGE achieves substantial FID improvements (15.8-25.9\%) while maintaining LPIPS within 7.0\% of baseline methods across both datasets.
    \item At $256 \times 256$ resolution on Flowers, FID improvements are smaller (9.1-15.0\%) but accompanied by larger LPIPS reductions (37.2-37.9\%), indicating reduced perceptual diversity at lower resolutions.
    \item LPIPS values decrease substantially from high to low resolution on Flowers (0.4209 to 0.1309), while the pattern differs across datasets based on their structural characteristics.
\end{itemize}